\section{Introduction}

\subsection{Motivation}

A neuron in a trained network rarely stands for one thing. It responds to
several unrelated inputs at once, a property usually called polysemanticity,
and it is the main obstacle to reading a model's computation off its weights.
One response to this is to change the architecture rather than the analysis:
if a layer is built out of many narrow experts, each expert may end up
responsible for a single feature, and the layer becomes legible by
construction.

\citet{park2025monet} push this idea furthest. Their MONET architecture stores
experts not individually but as compositions: a layer is cut into a
\emph{bottom} group of $\sqrt{N}$ factors and a \emph{top} group of
$\sqrt{N}$ factors, and expert $(i,j)$ is assembled on demand from bottom
factor $i$ and top factor $j$. The parameter count then grows as $\sqrt{N}$
rather than $N$, which lets them place $262\,144$ experts in a single layer
and report that experts specialise by domain: some route Python, others
chemistry, others toxic text.

The evidence for that specialisation is correlational. An expert is judged to
belong to a domain when its routing score is skewed towards that domain, and
toxicity experts are found by correlating routing scores with a toxicity
label. The authors are explicit that these criteria are rudimentary and that
quantitative evaluation of automated interpretability remains open. The deeper
difficulty is that on natural language there is no ground truth: nobody knows
what features a language model \emph{ought} to represent, so a claim that an
expert is monosemantic cannot be checked, only illustrated.

\subsection{Approach}

We therefore move the question to a task where the answer is known. For
$c = (a + b) \bmod P$ the generalising solution has been reverse-engineered
completely \citep{nanda2023progress}: the network represents each operand
through a small number of frequencies $\omega = 2\pi k / P$, computes
$\cos(\omega a + \varphi_1)$ and $\cos(\omega b + \varphi_2)$, and reads out
$\cos(\omega c + \varphi_1 + \varphi_2)$. The true features are Fourier
frequencies, and there are only $(P-1)/2$ of them. Monosemanticity becomes
measurable rather than merely illustrated: a unit is monosemantic exactly to the
extent that its response, viewed as a function of the operand, is a single
frequency.

The task carries a second advantage. Models trained on it exhibit grokking:
training accuracy saturates almost immediately while test accuracy stays at
chance for thousands of steps, and only then jumps. The step at which this
happens is a sharp, single-number measure of how hard it is for a given
architecture to find the generalising algorithm rather than a lookup table.

We replace the hidden layer of the reference MLP with seven alternatives ---
a dense baseline, an unfactorised MoE, MONET's horizontal and vertical
decompositions, and CP, Tucker and Tensor-Train factorisations of the expert
tensor in the style of \citet{oldfield2024mumoe} --- and sweep the expert
count from $16$ to $4096$ under an identical training protocol, with three
seeds per configuration.

\subsection{Contributions}

\begin{enumerate}
\item \textbf{The sign of the expert-count effect is set by the particular
  factorisation.} Adding experts brings generalisation forward for Tucker and
  CP layers, pushes it back for Tensor-Train and product-key layers, and
  destroys it outright for independent experts past $N = 64$. The claim that
  more experts is better is therefore not a property of mixtures of experts.
  Nor, contrary to what we expected at the outset, is it settled by whether
  the expert tensor is factorised at all: Tensor-Train shares a basis as
  thoroughly as Tucker and keeps it as clean over the middle of the range, yet
  scales in the opposite direction.

\item \textbf{A shared basis can act as a sparsity prior in the task's own
  basis, but does not always.} A Tucker-factorised layer solves modular
  addition using an effective $1.5$ to $1.8$ frequencies at every expert count
  we test, where a dense MLP spends $38.1$. The product-key layers begin in
  the same place and lose that compression as they grow, reaching $29.3$
  frequencies at $N = 4096$ with a spectral purity of $0.150$ against a
  white-noise floor of $0.082$. CP is the exception that fixes the scope of
  the claim: it shares a basis and broadens it, from $14.1$ to $35.6$
  frequencies, while its units become purer.

\item \textbf{Scaling drains the router of algorithmic structure, not of all
  structure.} As the expert count grows from $16$ to $4096$, the spectral
  purity of the router keys falls to the white-noise floor in every mixture
  architecture, at the largest expert count each is run at --- to
  $0.048$--$0.093$ against a floor of $0.082$ --- while a Tucker layer's
  shared basis holds purity above $0.95$. Yet on a two-subtask control in
  which the subtask is signalled in the input, the same routers at the same
  expert count support MONET's attribution criterion perfectly, selecting
  factors whose ablation destroys one subtask and spares the other. Routing
  scores recover partitions the input already carries and miss the structure
  that performs the computation, and nothing in the method distinguishes the
  two cases.

\item \textbf{A ground-truth testbed.} We build on the standard
  modular-addition setting a protocol in which claims about expert
  monosemanticity can be checked against a known answer, together with three
  metrics that measure different things --- effective frequency count,
  norm-weighted spectral purity, and router--basis spectral overlap --- and
  with code and run summaries released alongside.
\end{enumerate}

\subsection{Scope}

The title names transformers because the layers studied here are the
feed-forward blocks of transformer Mixture-of-Experts models, taken out of
that setting and trained in isolation. No transformer is trained in this work,
and the units we analyse are the factors of an expert layer rather than whole
experts. Two limitations bound these claims and are stated in full in
Section~\ref{sec:limitations}. The multilinear ($\mu$MoE) layers compute a
product of a function of $a$ and a function of $b$, which is structurally the
same form as the Fourier solution. Section~\ref{sec:res-ablations} measures how
much of their advantage follows from that alignment, by permuting the router's
input so that its two halves are no longer the two operands; the answer is that
a substantial part of it does. And all results are for a single task at a
single modulus, with a fixed train/test partition and three seeds varying
initialisation only, so they are reported as observations about this testbed
rather than as properties of mixture-of-experts layers in general.